\documentclass[11pt,reqno]{amsart}
\usepackage{graphicx}
\usepackage{indentfirst,csquotes}
\usepackage{etoolbox}
\usepackage{amssymb,amsthm,amsmath}
\usepackage{xcolor,paralist,hyperref,titlesec,fancyhdr,etoolbox}
% third highlight colour: dark violet (distinct from the blue/red terms)
\definecolor{darkviolet}{rgb}{0.45,0.10,0.62}
\usepackage{fancyhdr}
\usepackage{titlesec}
\usepackage{bbm}
\usepackage{etoolbox}
\usepackage{thmtools}
\usepackage{stmaryrd} % for \bigtimes
\usepackage{subcaption}
\usepackage{booktabs}
% \inputrows: table rows generated by issue174_cv_data.do and issue174_cv_did.py (../Data/*_rows.tex; make.sh), read inside a tabular. The TeX primitive
% \input is used because the LaTeX \input adds end-of-file code that breaks \bottomrule after the rows.
% If a rows file is not built yet, \inputrows prints a placeholder row instead of stopping.
\makeatletter\newcommand{\inputrows}[1]{\IfFileExists{#1}{\@@input #1 }{\gdef\cvhouseholds{??}\textit{Not yet built: \detokenize{#1}} & \\}}\makeatother
\providecommand{\cvhouseholds}{??} % households of the Table 3 regression; defined in ../Data/issue174_cv_did_gas_table3_rows.tex
\usepackage{float}
\usepackage{tikz}
\usepackage{pgfplots}
\pgfplotsset{compat=1.18}
\usepackage{pgfplots}
\pgfplotsset{compat=1.18}
\usetikzlibrary{decorations.markings}
\usetikzlibrary{arrows.meta}
\usepackage{hyperref}
\usepackage{marginnote}
\usepackage{enumitem}
\DeclareMathOperator*{\argmax}{arg\,max}
\usepackage[most]{tcolorbox}



%%%%%%%%%%%%%%%%%%%%%%% Vertical fill
% amsart is \flushbottom: it stretches the vertical glue so every page ends flush at the bottom
% margin. On a Section-4 page that carries only headings plus one tall [H] float, all of that
% stretch lands ABOVE the float -- the table drops to the bottom with a gap under the heading.
% \raggedbottom keeps each block at its natural position and lets the slack fall to the page foot.
\raggedbottom

%%%%%%%%%%%%%%%%%%%%%%% Float spacing
% Section 4 stacks a table and up to four figures per subsection, each its own [H] float. At the
% class defaults that reads as a column of isolated blocks separated by wide gaps. \intextsep is
% the one that matters for [H] (the float package places them in the text stream, not at the page
% top/bottom), but set all three so a float that does escape to a page break matches.
\setlength{\intextsep}{8pt plus 2pt minus 2pt}
\setlength{\textfloatsep}{8pt plus 2pt minus 2pt}
\setlength{\floatsep}{8pt plus 2pt minus 2pt}
\setlength{\abovecaptionskip}{4pt}
\setlength{\belowcaptionskip}{0pt}

%%%%%%%%%%%%%%%%%%%%%%% Margins
\usepackage[a4paper, top = 1cm, bottom = 1cm, left = 2cm, right = 2cm]{geometry}

%%%%%%%%%%%%%%%%%%%%%%% Environments
\newtheorem{theorem}{Theorem}[]
\newtheorem{definition}[theorem]{Definition}
\newtheorem{lemma}[theorem]{Lemma}
\newtheorem{proposition}[theorem]{Proposition}
\newtheorem{corollary}[theorem]{Corollary}
\newtheorem{conjecture}[theorem]{Conjecture}
\theoremstyle{definition}
\newtheorem{assumption}{Assumption}
\newtheorem{example}[theorem]{Example}
\newtheorem{remark}[theorem]{Remark}

%%%%%%%%%%%%%%%%%%%%%%% Footnotes
\hypersetup{colorlinks=true, linkcolor=blue, filecolor=black, urlcolor=black }

%%%%%%%%%%%%%%%%%%%%%%% Line-Spacing + Spacing Between Paragraphs
\linespread{1.0} 
\setlength{\parskip}{0.5em}
\newcommand{\Hquad}{\hspace{0.5em}} 

%%%%%%%%%%%%%%%%%%%%%%% Equation Spacing
\setlength{\abovedisplayskip}{6pt}  
\setlength{\belowdisplayskip}{6pt}
\setlength{\abovedisplayshortskip}{0pt}
\setlength{\belowdisplayshortskip}{4pt}

%%%%%%%%%%%%%%%%%%%%%%% Sections
\makeatletter

\providecommand{\@secnumpunct}{\quad} % Section title formatting
\renewcommand{\l@section}{\@tocline{1}{0pt}{1.5em}{2.3em}{}} % Table of Content Formatting
\renewcommand{\@secnumfont}{\bfseries} % Force amsart to display section numbers
\makeatother
\titleformat{\section}
  [block]
  {\normalfont\Large\bfseries}
  {\thesection}
  {1em}
  {}
\titleformat{\subsection}
  [block]
  {\normalfont\large\bfseries}
  {\thesubsection}
  {1em}
  {}

%%%%%%%%%%%%%%%%%%%%%%% Norms
\newcommand{\norm}[1]{\lVert #1 \rVert}

%%%%%%%%%%%%%%%%%%%%%%% Page Number Formatting
\pagestyle{fancy}
\fancyhf{}
\fancyfoot[C]{\thepage}  % Page number at bottom center
\renewcommand{\headrulewidth}{0pt}  % Remove header rule line
\setlength{\footskip}{13.0pt}

\begin{document}
\tableofcontents

\newpage
\section{Background}

\noindent In this document, I investigate if we observe success rate equalization across different committees within the same stream for the CA Insight Grants Program. 

\noindent There are 25 committees (equivalent to fields in our model) listed in the Table below. I split the committees into two groups: Group 1 and 2 containing 17 and 8 commitees respectively. In this document, I argue that there is evidence of (approximate) within stream success rate equalization across commitees \textit{within the same group} for competition years from 2020-2021 onwards. 


% The table is built by 2_analysis/2.2_descriptives/source/sshrc_figures.py; [H] keeps it here.
\floatplacement{table}{H}
\begin{table}
\centering
\small
% GENERATED BY sshrc_figures.py - edit that file, not this one.
% The document that inputs this file supplies the float, the caption and the label.
\begin{tabular}{ll}
\toprule
Code & Committee \\
\midrule
\multicolumn{2}{l}{\textit{Group 1: 17 committees funded at a common rate within each stream}} \\
\addlinespace[2pt]
01 & Philosophy \\
02 & History \\
03 & Fine Arts, Research-Creation \\
04 & Literature \\
05 & Medieval, classics, religious studies \\
08 & Sociology, demography and related fields \\
09 & Geography, urban planning and related fields \\
10A & Psychology \\
10B & Psychology \\
11 & Political science and public administration \\
12A & Education \\
12B & Education and social work \\
12C & Education and social work \\
13 & Anthropology and archaeology \\
15 & Linguistics and translation \\
16 & Communications, media studies, gender studies, library and information science, related fields \\
17 & Law and criminology \\
\midrule
\multicolumn{2}{l}{\textit{Group 2: the other 8 committees}} \\
\addlinespace[2pt]
07 & Economics \\
14A & Business, management and related fields \\
14B & Business, management and related fields \\
14C & Business, management and related fields \\
22 & Multidisciplinary or interdisciplinary humanities \\
23 & Multidisciplinary or interdisciplinary social sciences \\
23A & Multidisciplinary or interdisciplinary social sciences \\
23B & Multidisciplinary or interdisciplinary social sciences \\
\bottomrule
\end{tabular}

\caption{\footnotesize SSHRC Insight Grant merit review committees by group, with Multidisciplinary social sciences (23) split into 23A and 23B from 2022--23 and Indigenous research (21) and the Tri-Agency interdisciplinary committee (24) in neither group.}
\label{tab:sshrc_committees}
\end{table}

\newpage
\floatplacement{figure}{H}

\section{Success Rates across time}

\noindent Graphs below plots the success rate of the commitees in each group $g\in \{1,2\}$ and stream $A,B$ combination, for all years after the reform of the Insight Grants program in 2017. The solid line is the pooled success rate (total awards / total eligible applicants). For each Figure, we observe the following two items for competition years from 2020-2021 onwards:
\begin{enumerate}[label=(\arabic*),ref=\thelemma(\arabic*)]
  \item Within each competition year, there is little dispersion in success rates across different commitees. 
  \item Interestingly, the commitee-specific success rates track each other over time. For example, consider Figure \ref{fig:sshrc_rates_common_A}. There are some years where the pooled success rate is high (e.g., 2022-2023) and where the pooled success rate is low (e.g. 2024-2025). What is interesting is that \textit{all} the individual commitee specific success rates also exhibit the same temporal trend: The individual success rates are all large in 2022-2023, and all low in 2024-2025 (i.e., they seem to track one another). 
\end{enumerate}

\begin{figure}
\centering
% GENERATED BY sshrc_figures.py - edit that file, not this one.
% The document that inputs this file supplies the float, the caption and the label.
% Stream A. Committees, left to right within each year: Philosophy, History, Fine arts, Literature, Sociology, Geography, Psychology (10A, 10B), Political science, Education (12A, 12B, 12C), Anthropology, Linguistics, Communications, Law and criminology.
\definecolor{sshrcA}{HTML}{2A78D6}
\definecolor{sshrcB}{HTML}{EB6834}
\begin{tikzpicture}
\begin{axis}[
  width=15cm, height=5.5cm,
  xmin=2017.6, xmax=2025.4, ymin=20, ymax=70,
  xtick={2018,2019,2020,2021,2022,2023,2024,2025}, xticklabels={2018--19,2019--20,2020--21,2021--22,2022--23,2023--24,2024--25,2025--26},
  axis x line=bottom, axis y line=left, axis line style={gray!70},
  ymajorgrids, grid style={gray!20},
  tick label style={font=\footnotesize}, label style={font=\footnotesize},
  xlabel={SSHRC results year}, ylabel={Success rate}, ylabel style={align=center},
  clip=false,
  ytick distance=10, yticklabel={\pgfmathprintnumber{\tick}\%},
]
\draw[gray, dashed, line width=0.7pt] (axis cs:2020,20) -- (axis cs:2020,70);
\addplot[only marks, mark=*, mark size=1.8pt, color=sshrcA,
  mark options={fill opacity=0.6, draw opacity=0}]
  coordinates {(2018.18,52.94) (2018.26,42.42) (2018.06,51.85) (2018.10,50.00) (2018.14,52.00) (2017.78,55.56) (2017.90,54.55) (2017.74,43.33) (2018.30,48.57) (2018.22,50.00) (2017.82,57.14) (2017.70,47.50) (2018.02,57.14) (2017.94,37.93) (2017.98,53.12) (2017.86,54.29) (2019.18,55.56) (2019.26,48.15) (2019.06,46.88) (2019.10,45.71) (2018.78,50.00) (2018.90,47.37) (2018.74,48.84) (2019.30,48.00) (2019.22,50.00) (2018.82,52.17) (2018.70,50.00) (2019.02,50.00) (2018.94,38.89) (2018.98,47.37) (2018.86,48.72) (2019.70,46.51) (2019.74,43.64) (2019.78,42.86) (2019.82,44.64) (2019.86,42.31) (2019.90,42.86) (2019.94,44.00) (2019.98,45.83) (2020.02,45.71) (2020.06,43.48) (2020.10,42.31) (2020.18,50.00) (2020.22,46.67) (2020.26,46.15) (2020.30,44.44) (2020.70,56.67) (2020.74,55.88) (2020.78,55.88) (2020.82,53.12) (2020.86,54.55) (2020.90,55.56) (2020.94,55.00) (2020.98,55.56) (2021.02,54.84) (2021.06,56.25) (2021.10,53.33) (2021.14,57.14) (2021.18,50.00) (2021.22,54.55) (2021.26,56.52) (2021.30,61.90) (2021.70,53.85) (2021.74,57.89) (2021.78,55.56) (2021.82,53.57) (2021.86,57.89) (2021.90,57.14) (2021.94,53.85) (2021.98,54.55) (2022.02,56.52) (2022.06,54.55) (2022.10,53.85) (2022.14,58.33) (2022.18,50.00) (2022.22,50.00) (2022.26,52.63) (2022.30,52.00) (2022.70,46.15) (2022.74,50.00) (2022.78,64.29) (2022.82,48.15) (2022.86,55.56) (2022.90,45.45) (2022.94,46.15) (2022.98,47.06) (2023.02,45.83) (2023.06,50.00) (2023.10,50.00) (2023.14,50.00) (2023.18,33.33) (2023.22,50.00) (2023.26,50.00) (2023.30,50.00) (2023.70,38.89) (2023.74,35.48) (2023.78,35.00) (2023.82,39.13) (2023.86,37.50) (2023.90,33.33) (2023.94,40.00) (2023.98,41.18) (2024.02,36.67) (2024.06,40.00) (2024.10,36.36) (2024.14,35.71) (2024.18,40.00) (2024.22,33.33) (2024.26,40.74) (2024.30,42.11) (2024.70,41.18) (2024.74,42.11) (2024.78,40.91) (2024.82,41.67) (2024.86,42.86) (2024.90,46.15) (2024.94,43.75) (2024.98,44.44) (2025.02,43.24) (2025.06,43.48) (2025.10,45.45) (2025.14,44.44) (2025.18,41.67) (2025.22,42.86) (2025.26,42.31) (2025.30,42.31)};
\addplot[color=sshrcA, line width=1.2pt, mark=*, mark size=2pt] coordinates {(2018.00,50.10) (2019.00,48.59) (2020.00,44.58) (2021.00,55.65) (2022.00,54.51) (2023.00,49.24) (2024.00,38.03) (2025.00,42.86)};
\end{axis}
\end{tikzpicture}

\caption{\footnotesize Group 1, Stream~A}
\label{fig:sshrc_rates_common_A}
\end{figure}

\begin{figure}
\centering
% GENERATED BY sshrc_figures.py - edit that file, not this one.
% The document that inputs this file supplies the float, the caption and the label.
% Stream B. Committees, left to right within each year: Philosophy, History, Fine arts, Literature, Medieval, classics and religious studies, Sociology, Geography, Psychology (10A, 10B), Political science, Education (12A, 12B, 12C), Anthropology, Linguistics, Communications, Law and criminology.
\definecolor{sshrcA}{HTML}{2A78D6}
\definecolor{sshrcB}{HTML}{EB6834}
\begin{tikzpicture}
\begin{axis}[
  width=15cm, height=5.5cm,
  xmin=2017.6, xmax=2025.4, ymin=20, ymax=70,
  xtick={2018,2019,2020,2021,2022,2023,2024,2025}, xticklabels={2018--19,2019--20,2020--21,2021--22,2022--23,2023--24,2024--25,2025--26},
  axis x line=bottom, axis y line=left, axis line style={gray!70},
  ymajorgrids, grid style={gray!20},
  tick label style={font=\footnotesize}, label style={font=\footnotesize},
  xlabel={SSHRC results year}, ylabel={Success rate}, ylabel style={align=center},
  clip=false,
  ytick distance=10, yticklabel={\pgfmathprintnumber{\tick}\%},
]
\draw[gray, dashed, line width=0.7pt] (axis cs:2020,20) -- (axis cs:2020,70);
\addplot[only marks, mark=square*, mark size=1.8pt, color=sshrcB,
  mark options={fill opacity=0.6, draw opacity=0}]
  coordinates {(2018.06,44.12) (2018.09,33.33) (2018.03,43.48) (2018.04,40.54) (2018.05,43.75) (2017.92,46.81) (2017.97,44.90) (2017.91,34.62) (2018.10,39.47) (2018.08,43.48) (2017.94,47.83) (2017.95,47.06) (2017.90,37.50) (2018.01,47.73) (2017.99,31.25) (2018.00,44.44) (2017.96,45.00) (2019.06,44.12) (2019.09,40.62) (2019.03,38.46) (2019.04,35.71) (2018.92,40.82) (2018.97,36.36) (2018.91,39.39) (2019.10,38.24) (2019.08,44.44) (2018.94,44.44) (2018.95,41.67) (2018.90,42.86) (2019.01,54.29) (2018.99,39.53) (2019.00,37.78) (2018.96,40.00) (2019.90,37.50) (2019.91,33.33) (2019.92,34.04) (2019.94,35.71) (2019.95,41.67) (2019.96,34.00) (2019.97,32.43) (2019.99,35.42) (2020.00,36.73) (2020.01,35.19) (2020.03,35.29) (2020.04,34.85) (2020.06,39.39) (2020.08,36.36) (2020.09,41.18) (2020.10,35.90) (2020.90,33.33) (2020.91,52.63) (2020.92,48.89) (2020.94,40.00) (2020.95,50.00) (2020.96,50.00) (2020.97,48.65) (2020.99,52.38) (2021.00,50.00) (2021.01,48.78) (2021.03,43.59) (2021.04,47.50) (2021.05,48.78) (2021.06,48.00) (2021.08,46.67) (2021.09,50.00) (2021.10,50.00) (2021.90,33.33) (2021.91,47.62) (2021.92,48.57) (2021.94,46.15) (2021.95,53.33) (2021.96,46.67) (2021.97,50.00) (2021.99,44.44) (2022.00,45.24) (2022.01,48.39) (2022.03,47.22) (2022.04,47.06) (2022.05,48.57) (2022.06,42.86) (2022.08,42.86) (2022.09,45.95) (2022.10,45.45) (2022.90,40.00) (2022.91,41.18) (2022.92,41.18) (2022.94,43.75) (2022.95,44.44) (2022.96,41.67) (2022.97,41.46) (2022.99,37.78) (2023.00,37.78) (2023.01,37.14) (2023.03,38.89) (2023.04,37.84) (2023.05,40.54) (2023.06,36.00) (2023.08,40.00) (2023.09,41.38) (2023.10,42.31) (2023.90,33.33) (2023.91,25.00) (2023.92,26.92) (2023.94,31.82) (2023.95,33.33) (2023.96,27.91) (2023.97,26.00) (2023.99,29.63) (2024.00,27.87) (2024.01,27.87) (2024.03,31.91) (2024.04,30.00) (2024.05,29.79) (2024.06,26.32) (2024.08,32.14) (2024.09,26.42) (2024.10,25.81) (2024.90,36.36) (2024.91,35.71) (2024.92,35.21) (2024.94,34.48) (2024.95,36.36) (2024.96,34.00) (2024.97,36.67) (2024.99,36.62) (2025.00,36.76) (2025.01,34.72) (2025.03,36.00) (2025.04,33.33) (2025.05,34.00) (2025.06,34.88) (2025.08,36.00) (2025.09,34.78) (2025.10,36.67)};
\addplot[color=sshrcB, line width=1.2pt, mark=square*, mark size=2pt] coordinates {(2018.00,42.52) (2019.00,40.34) (2020.00,35.70) (2021.00,48.63) (2022.00,46.71) (2023.00,39.84) (2024.00,28.53) (2025.00,35.38)};
\end{axis}
\end{tikzpicture}

\caption{\footnotesize Group 1, Stream~B}
\label{fig:sshrc_rates_common_B}
\end{figure}


\begin{figure}
\centering
% GENERATED BY sshrc_figures.py - edit that file, not this one.
% The document that inputs this file supplies the float, the caption and the label.
% Stream A. Committees, left to right within each year: Medieval, classics and religious studies, Economics, Business (14A, 14B, 14C), Multidisciplinary humanities, Multidisciplinary social sciences (23, 23A, 23B).
\definecolor{sshrcA}{HTML}{2A78D6}
\definecolor{sshrcB}{HTML}{EB6834}
\begin{tikzpicture}
\begin{axis}[
  width=15cm, height=5.5cm,
  xmin=2017.6, xmax=2025.4, ymin=20, ymax=70,
  xtick={2018,2019,2020,2021,2022,2023,2024,2025}, xticklabels={2018--19,2019--20,2020--21,2021--22,2022--23,2023--24,2024--25,2025--26},
  axis x line=bottom, axis y line=left, axis line style={gray!70},
  ymajorgrids, grid style={gray!20},
  tick label style={font=\footnotesize}, label style={font=\footnotesize},
  xlabel={SSHRC results year}, ylabel={Success rate}, ylabel style={align=center},
  clip=false,
  ytick distance=10, yticklabel={\pgfmathprintnumber{\tick}\%},
]
\draw[gray, dashed, line width=0.7pt] (axis cs:2020,20) -- (axis cs:2020,70);
\addplot[only marks, mark=*, mark size=1.8pt, color=sshrcA,
  mark options={fill opacity=0.6, draw opacity=0}]
  coordinates {(2017.85,60.00) (2017.92,62.50) (2017.78,57.38) (2017.70,55.00) (2018.08,40.91) (2018.15,40.00) (2018.85,58.97) (2018.92,57.50) (2018.78,57.78) (2018.70,52.63) (2019.08,53.33) (2019.15,51.85) (2019.70,45.00) (2019.78,46.34) (2019.85,47.50) (2019.92,48.57) (2020.08,47.83) (2020.15,48.15) (2020.70,55.00) (2020.78,58.06) (2020.85,59.09) (2020.92,58.33) (2021.00,62.50) (2021.08,57.14) (2021.15,58.33) (2021.70,53.85) (2021.78,60.61) (2021.85,58.33) (2021.92,60.00) (2022.00,61.54) (2022.08,61.54) (2022.22,66.67) (2022.30,60.00) (2022.70,50.00) (2022.78,51.22) (2022.85,54.17) (2022.92,52.17) (2023.00,50.00) (2023.08,60.00) (2023.22,42.86) (2023.30,66.67) (2023.70,47.37) (2023.78,43.64) (2023.85,42.86) (2023.92,43.33) (2024.00,44.44) (2024.08,46.15) (2024.22,46.15) (2024.30,50.00) (2024.70,44.00) (2024.78,49.30) (2024.85,48.78) (2024.92,47.37) (2025.00,47.50) (2025.08,50.00) (2025.22,58.33) (2025.30,54.55)};
\addplot[color=sshrcA, line width=1.2pt, mark=*, mark size=2pt] coordinates {(2018.00,54.71) (2019.00,56.22) (2020.00,47.31) (2021.00,58.55) (2022.00,60.13) (2023.00,52.41) (2024.00,44.67) (2025.00,48.77)};
\end{axis}
\end{tikzpicture}

\caption{\footnotesize Group 2, Stream~A}
\label{fig:sshrc_rates_above_A}
\end{figure}

\begin{figure}
\centering
% GENERATED BY sshrc_figures.py - edit that file, not this one.
% The document that inputs this file supplies the float, the caption and the label.
% Stream B. Committees, left to right within each year: Economics, Business (14A, 14B, 14C), Multidisciplinary humanities, Multidisciplinary social sciences (23, 23A, 23B).
\definecolor{sshrcA}{HTML}{2A78D6}
\definecolor{sshrcB}{HTML}{EB6834}
\begin{tikzpicture}
\begin{axis}[
  width=15cm, height=5.5cm,
  xmin=2017.6, xmax=2025.4, ymin=20, ymax=70,
  xtick={2018,2019,2020,2021,2022,2023,2024,2025}, xticklabels={2018--19,2019--20,2020--21,2021--22,2022--23,2023--24,2024--25,2025--26},
  axis x line=bottom, axis y line=left, axis line style={gray!70},
  ymajorgrids, grid style={gray!20},
  tick label style={font=\footnotesize}, label style={font=\footnotesize},
  xlabel={SSHRC results year}, ylabel={Success rate}, ylabel style={align=center},
  clip=false,
  ytick distance=10, yticklabel={\pgfmathprintnumber{\tick}\%},
]
\draw[gray, dashed, line width=0.7pt] (axis cs:2020,20) -- (axis cs:2020,70);
\addplot[only marks, mark=square*, mark size=1.8pt, color=sshrcB,
  mark options={fill opacity=0.6, draw opacity=0}]
  coordinates {(2017.93,51.11) (2017.96,52.94) (2017.90,48.00) (2018.01,38.89) (2018.04,38.98) (2018.93,48.57) (2018.96,48.57) (2018.90,47.83) (2019.01,42.86) (2019.04,43.48) (2019.90,37.93) (2019.93,42.50) (2019.96,43.18) (2020.01,37.50) (2020.04,43.10) (2020.90,56.00) (2020.93,51.85) (2020.96,52.00) (2020.99,56.00) (2021.01,50.00) (2021.04,49.32) (2021.90,50.00) (2021.93,51.85) (2021.96,51.85) (2021.99,52.00) (2022.01,50.00) (2022.07,54.84) (2022.10,55.56) (2022.90,42.86) (2022.93,44.00) (2022.96,42.86) (2022.99,42.31) (2023.01,47.37) (2023.07,45.45) (2023.10,46.88) (2023.90,35.29) (2023.93,35.00) (2023.96,33.33) (2023.99,33.33) (2024.01,38.89) (2024.07,36.36) (2024.10,36.36) (2024.90,42.42) (2024.93,41.46) (2024.96,41.86) (2024.99,40.48) (2025.01,45.00) (2025.07,45.71) (2025.10,45.45)};
\addplot[color=sshrcB, line width=1.2pt, mark=square*, mark size=2pt] coordinates {(2018.00,45.86) (2019.00,46.41) (2020.00,41.54) (2021.00,51.83) (2022.00,52.63) (2023.00,44.50) (2024.00,35.17) (2025.00,42.91)};
\end{axis}
\end{tikzpicture}

\caption{\footnotesize Group 2, Stream~B}
\label{fig:sshrc_rates_above_B}
\end{figure}

\noindent In the two figures below, I plot the success rates of committees in group 1 with those of group 2 for each stream. In general, within stream success rates of committees in group 2 are larger than the success rates of committees in group 1. 

\begin{figure}
\centering
% GENERATED BY sshrc_figures.py - edit that file, not this one.
% The document that inputs this file supplies the float, the caption and the label.
% Stream A. Group 1: Philosophy, History, Fine arts, Literature, Medieval, classics and religious studies, Sociology, Geography, Psychology (10A, 10B), Political science, Education (12A, 12B, 12C), Anthropology, Linguistics, Communications, Law and criminology. Group 2: Economics, Business (14A, 14B, 14C), Multidisciplinary humanities, Multidisciplinary social sciences (23, 23A, 23B).
\definecolor{sshrccommon}{HTML}{1BAF7A}
\definecolor{sshrcabove}{HTML}{6250D6}
\begin{tikzpicture}
\begin{axis}[
  width=15cm, height=5.5cm,
  xmin=2017.6, xmax=2025.4, ymin=20, ymax=70,
  xtick={2018,2019,2020,2021,2022,2023,2024,2025}, xticklabels={2018--19,2019--20,2020--21,2021--22,2022--23,2023--24,2024--25,2025--26},
  axis x line=bottom, axis y line=left, axis line style={gray!70},
  ymajorgrids, grid style={gray!20},
  tick label style={font=\footnotesize}, label style={font=\footnotesize},
  xlabel={SSHRC results year}, ylabel={Success rate}, ylabel style={align=center},
  clip=false,
  ytick distance=10, yticklabel={\pgfmathprintnumber{\tick}\%},
  legend style={font=\footnotesize, draw=none, fill=none, at={(0.5,1.02)}, anchor=south,
    legend columns=2, /tikz/every even column/.append style={column sep=1em}},
  legend cell align=left,
]
\draw[gray, dashed, line width=0.7pt] (axis cs:2020,20) -- (axis cs:2020,70);
\addplot[only marks, mark=*, mark size=1.8pt, color=sshrccommon,
  mark options={fill opacity=0.5, draw opacity=0}]
  coordinates {(2018.01,52.94) (2018.03,42.42) (2017.98,51.85) (2017.99,50.00) (2018.00,52.00) (2017.92,55.56) (2017.95,54.55) (2017.91,43.33) (2018.03,48.57) (2018.02,50.00) (2017.92,57.14) (2017.93,55.00) (2017.90,47.50) (2017.97,57.14) (2017.96,37.93) (2017.97,53.12) (2017.94,54.29) (2019.01,55.56) (2019.03,48.15) (2018.98,46.88) (2018.99,45.71) (2018.92,50.00) (2018.95,47.37) (2018.91,48.84) (2019.03,48.00) (2019.02,50.00) (2018.92,52.17) (2018.93,52.63) (2018.90,50.00) (2018.97,50.00) (2018.96,38.89) (2018.97,47.37) (2018.94,48.72) (2019.90,46.51) (2019.91,43.64) (2019.92,42.86) (2019.92,44.64) (2019.93,45.00) (2019.94,42.31) (2019.95,42.86) (2019.96,44.00) (2019.97,45.83) (2019.97,45.71) (2019.98,43.48) (2019.99,42.31) (2020.01,50.00) (2020.02,46.67) (2020.03,46.15) (2020.03,44.44) (2020.90,56.67) (2020.91,55.88) (2020.92,55.88) (2020.92,53.12) (2020.93,55.00) (2020.94,54.55) (2020.95,55.56) (2020.96,55.00) (2020.97,55.56) (2020.97,54.84) (2020.98,56.25) (2020.99,53.33) (2021.00,57.14) (2021.01,50.00) (2021.02,54.55) (2021.03,56.52) (2021.03,61.90) (2021.90,53.85) (2021.91,57.89) (2021.92,55.56) (2021.92,53.57) (2021.93,53.85) (2021.94,57.89) (2021.95,57.14) (2021.96,53.85) (2021.97,54.55) (2021.97,56.52) (2021.98,54.55) (2021.99,53.85) (2022.00,58.33) (2022.01,50.00) (2022.02,50.00) (2022.03,52.63) (2022.03,52.00) (2022.90,46.15) (2022.91,50.00) (2022.92,64.29) (2022.92,48.15) (2022.93,50.00) (2022.94,55.56) (2022.95,45.45) (2022.96,46.15) (2022.97,47.06) (2022.97,45.83) (2022.98,50.00) (2022.99,50.00) (2023.00,50.00) (2023.01,33.33) (2023.02,50.00) (2023.03,50.00) (2023.03,50.00) (2023.90,38.89) (2023.91,35.48) (2023.92,35.00) (2023.92,39.13) (2023.93,47.37) (2023.94,37.50) (2023.95,33.33) (2023.96,40.00) (2023.97,41.18) (2023.97,36.67) (2023.98,40.00) (2023.99,36.36) (2024.00,35.71) (2024.01,40.00) (2024.02,33.33) (2024.03,40.74) (2024.03,42.11) (2024.90,41.18) (2024.91,42.11) (2024.92,40.91) (2024.92,41.67) (2024.93,44.00) (2024.94,42.86) (2024.95,46.15) (2024.96,43.75) (2024.97,44.44) (2024.97,43.24) (2024.98,43.48) (2024.99,45.45) (2025.00,44.44) (2025.01,41.67) (2025.02,42.86) (2025.03,42.31) (2025.03,42.31)};
\addlegendentry{Group 1 (17 committees)}
\addplot[only marks, mark=triangle*, mark size=2.4pt, color=sshrcabove,
  mark options={fill opacity=0.5, draw opacity=0}]
  coordinates {(2018.05,60.00) (2018.06,62.50) (2018.04,57.38) (2018.08,40.91) (2018.08,40.00) (2019.05,58.97) (2019.06,57.50) (2019.04,57.78) (2019.08,53.33) (2019.08,51.85) (2020.04,46.34) (2020.05,47.50) (2020.06,48.57) (2020.08,47.83) (2020.08,48.15) (2021.04,58.06) (2021.05,59.09) (2021.06,58.33) (2021.07,62.50) (2021.08,57.14) (2021.08,58.33) (2022.04,60.61) (2022.05,58.33) (2022.06,60.00) (2022.07,61.54) (2022.08,61.54) (2022.09,66.67) (2022.10,60.00) (2023.04,51.22) (2023.05,54.17) (2023.06,52.17) (2023.07,50.00) (2023.08,60.00) (2023.09,42.86) (2023.10,66.67) (2024.04,43.64) (2024.05,42.86) (2024.06,43.33) (2024.07,44.44) (2024.08,46.15) (2024.09,46.15) (2024.10,50.00) (2025.04,49.30) (2025.05,48.78) (2025.06,47.37) (2025.07,47.50) (2025.08,50.00) (2025.09,58.33) (2025.10,54.55)};
\addlegendentry{Group 2 (8 committees)}
\end{axis}
\end{tikzpicture}

\caption{\footnotesize Groups 1 and 2, Stream~A}
\label{fig:sshrc_rates_both_A}
\end{figure}

\begin{figure}
\centering
% GENERATED BY sshrc_figures.py - edit that file, not this one.
% The document that inputs this file supplies the float, the caption and the label.
% Stream B. Group 1: Philosophy, History, Fine arts, Literature, Medieval, classics and religious studies, Sociology, Geography, Psychology (10A, 10B), Political science, Education (12A, 12B, 12C), Anthropology, Linguistics, Communications, Law and criminology. Group 2: Economics, Business (14A, 14B, 14C), Multidisciplinary humanities, Multidisciplinary social sciences (23, 23A, 23B).
\definecolor{sshrccommon}{HTML}{1BAF7A}
\definecolor{sshrcabove}{HTML}{6250D6}
\begin{tikzpicture}
\begin{axis}[
  width=15cm, height=5.5cm,
  xmin=2017.6, xmax=2025.4, ymin=20, ymax=70,
  xtick={2018,2019,2020,2021,2022,2023,2024,2025}, xticklabels={2018--19,2019--20,2020--21,2021--22,2022--23,2023--24,2024--25,2025--26},
  axis x line=bottom, axis y line=left, axis line style={gray!70},
  ymajorgrids, grid style={gray!20},
  tick label style={font=\footnotesize}, label style={font=\footnotesize},
  xlabel={SSHRC results year}, ylabel={Success rate}, ylabel style={align=center},
  clip=false,
  ytick distance=10, yticklabel={\pgfmathprintnumber{\tick}\%},
  legend style={font=\footnotesize, draw=none, fill=none, at={(0.5,1.02)}, anchor=south,
    legend columns=2, /tikz/every even column/.append style={column sep=1em}},
  legend cell align=left,
]
\draw[gray, dashed, line width=0.7pt] (axis cs:2020,20) -- (axis cs:2020,70);
\addplot[only marks, mark=*, mark size=1.8pt, color=sshrccommon,
  mark options={fill opacity=0.5, draw opacity=0}]
  coordinates {(2018.01,44.12) (2018.03,33.33) (2017.98,43.48) (2017.99,40.54) (2018.00,43.75) (2017.92,46.81) (2017.95,44.90) (2017.91,34.62) (2018.03,39.47) (2018.02,43.48) (2017.92,47.83) (2017.93,47.06) (2017.90,37.50) (2017.97,47.73) (2017.96,31.25) (2017.97,44.44) (2017.94,45.00) (2019.01,44.12) (2019.03,40.62) (2018.98,38.46) (2018.99,35.71) (2018.92,40.82) (2018.95,36.36) (2018.91,39.39) (2019.03,38.24) (2019.02,44.44) (2018.92,44.44) (2018.93,41.67) (2018.90,42.86) (2018.97,54.29) (2018.96,39.53) (2018.97,37.78) (2018.94,40.00) (2019.90,37.50) (2019.91,33.33) (2019.92,34.04) (2019.92,35.71) (2019.93,41.67) (2019.94,34.00) (2019.95,32.43) (2019.96,35.42) (2019.97,36.73) (2019.97,35.19) (2019.98,35.29) (2019.99,34.85) (2020.01,39.39) (2020.02,36.36) (2020.03,41.18) (2020.03,35.90) (2020.90,33.33) (2020.91,52.63) (2020.92,48.89) (2020.92,40.00) (2020.93,50.00) (2020.94,50.00) (2020.95,48.65) (2020.96,52.38) (2020.97,50.00) (2020.97,48.78) (2020.98,43.59) (2020.99,47.50) (2021.00,48.78) (2021.01,48.00) (2021.02,46.67) (2021.03,50.00) (2021.03,50.00) (2021.90,33.33) (2021.91,47.62) (2021.92,48.57) (2021.92,46.15) (2021.93,53.33) (2021.94,46.67) (2021.95,50.00) (2021.96,44.44) (2021.97,45.24) (2021.97,48.39) (2021.98,47.22) (2021.99,47.06) (2022.00,48.57) (2022.01,42.86) (2022.02,42.86) (2022.03,45.95) (2022.03,45.45) (2022.90,40.00) (2022.91,41.18) (2022.92,41.18) (2022.92,43.75) (2022.93,44.44) (2022.94,41.67) (2022.95,41.46) (2022.96,37.78) (2022.97,37.78) (2022.97,37.14) (2022.98,38.89) (2022.99,37.84) (2023.00,40.54) (2023.01,36.00) (2023.02,40.00) (2023.03,41.38) (2023.03,42.31) (2023.90,33.33) (2023.91,25.00) (2023.92,26.92) (2023.92,31.82) (2023.93,33.33) (2023.94,27.91) (2023.95,26.00) (2023.96,29.63) (2023.97,27.87) (2023.97,27.87) (2023.98,31.91) (2023.99,30.00) (2024.00,29.79) (2024.01,26.32) (2024.02,32.14) (2024.03,26.42) (2024.03,25.81) (2024.90,36.36) (2024.91,35.71) (2024.92,35.21) (2024.92,34.48) (2024.93,36.36) (2024.94,34.00) (2024.95,36.67) (2024.96,36.62) (2024.97,36.76) (2024.97,34.72) (2024.98,36.00) (2024.99,33.33) (2025.00,34.00) (2025.01,34.88) (2025.02,36.00) (2025.03,34.78) (2025.03,36.67)};
\addlegendentry{Group 1 (17 committees)}
\addplot[only marks, mark=triangle*, mark size=2.4pt, color=sshrcabove,
  mark options={fill opacity=0.5, draw opacity=0}]
  coordinates {(2018.05,51.11) (2018.06,52.94) (2018.04,48.00) (2018.08,38.89) (2018.08,38.98) (2019.05,48.57) (2019.06,48.57) (2019.04,47.83) (2019.08,42.86) (2019.08,43.48) (2020.04,37.93) (2020.05,42.50) (2020.06,43.18) (2020.08,37.50) (2020.08,43.10) (2021.04,56.00) (2021.05,51.85) (2021.06,52.00) (2021.07,56.00) (2021.08,50.00) (2021.08,49.32) (2022.04,50.00) (2022.05,51.85) (2022.06,51.85) (2022.07,52.00) (2022.08,50.00) (2022.09,54.84) (2022.10,55.56) (2023.04,42.86) (2023.05,44.00) (2023.06,42.86) (2023.07,42.31) (2023.08,47.37) (2023.09,45.45) (2023.10,46.88) (2024.04,35.29) (2024.05,35.00) (2024.06,33.33) (2024.07,33.33) (2024.08,38.89) (2024.09,36.36) (2024.10,36.36) (2025.04,42.42) (2025.05,41.46) (2025.06,41.86) (2025.07,40.48) (2025.08,45.00) (2025.09,45.71) (2025.10,45.45)};
\addlegendentry{Group 2 (8 committees)}
\end{axis}
\end{tikzpicture}

\caption{\footnotesize Groups 1 and 2, Stream~B}
\label{fig:sshrc_rates_both_B}
\end{figure}

\noindent 


% Spread figure, dropped for the moment:
% \begin{figure}
% \centering
% % GENERATED BY sshrc_figures.py - edit that file, not this one.
% The document that inputs this file supplies the float, the caption and the label.
% Stream A. Committees: Philosophy, History, Fine arts, Literature, Medieval, classics and religious studies, Sociology, Geography, Psychology (10A, 10B), Political science, Education (12A, 12B, 12C), Anthropology, Linguistics, Communications, Law and criminology.
% spread (Pearson chi2/df) and its whole-award floor: 2018: 0.402 (0.015); 2019: 0.099 (0.016); 2020: 0.027 (0.018); 2021: 0.037 (0.013); 2022: 0.038 (0.023); 2023: 0.132 (0.037); 2024: 0.091 (0.027); 2025: 0.019 (0.016)
\definecolor{sshrcA}{HTML}{2A78D6}
\definecolor{sshrcB}{HTML}{EB6834}
\begin{tikzpicture}
\begin{axis}[
  width=15cm, height=5.5cm,
  xmin=2017.6, xmax=2025.4, ymin=0, ymax=1.75,
  xtick={2018,2019,2020,2021,2022,2023,2024,2025}, xticklabels={2018--19,2019--20,2020--21,2021--22,2022--23,2023--24,2024--25,2025--26},
  axis x line=bottom, axis y line=left, axis line style={gray!70},
  ymajorgrids, grid style={gray!20},
  tick label style={font=\footnotesize}, label style={font=\footnotesize},
  xlabel={SSHRC results year}, ylabel={Spread of committee success rates}, ylabel style={align=center},
  clip=false,
  ytick distance=0.25, yticklabel style={/pgf/number format/.cd, fixed, fixed zerofill, precision=2},
]
\draw[gray, dashed, line width=0.7pt] (axis cs:2020,0) -- (axis cs:2020,1.75);
\addplot[gray, line width=0.6pt, densely dotted, forget plot, domain=2017.6:2025.4, samples=2] {1};
\node[font=\footnotesize, text=gray!80!black, anchor=south east] at (axis cs:2025.4,1) {chance};
\addplot[color=sshrcA, line width=1.2pt, mark=*, mark size=1.8pt] coordinates {(2018.00,0.40) (2019.00,0.10) (2020.00,0.03) (2021.00,0.04) (2022.00,0.04) (2023.00,0.13) (2024.00,0.09) (2025.00,0.02)};
\end{axis}
\end{tikzpicture}

% \caption{\footnotesize Group 1, Stream~A}
% \label{fig:sshrc_spread_common_A}
% \end{figure}



% Spread figure, dropped for the moment:
% \begin{figure}
% \centering
% % GENERATED BY sshrc_figures.py - edit that file, not this one.
% The document that inputs this file supplies the float, the caption and the label.
% Stream B. Committees: Philosophy, History, Fine arts, Literature, Medieval, classics and religious studies, Sociology, Geography, Psychology (10A, 10B), Political science, Education (12A, 12B, 12C), Anthropology, Linguistics, Communications, Law and criminology.
% spread (Pearson chi2/df) and its whole-award floor: 2018: 0.343 (0.014); 2019: 0.305 (0.007); 2020: 0.083 (0.010); 2021: 0.103 (0.030); 2022: 0.074 (0.021); 2023: 0.062 (0.012); 2024: 0.108 (0.016); 2025: 0.026 (0.007)
\definecolor{sshrcA}{HTML}{2A78D6}
\definecolor{sshrcB}{HTML}{EB6834}
\begin{tikzpicture}
\begin{axis}[
  width=15cm, height=5.5cm,
  xmin=2017.6, xmax=2025.4, ymin=0, ymax=1.75,
  xtick={2018,2019,2020,2021,2022,2023,2024,2025}, xticklabels={2018--19,2019--20,2020--21,2021--22,2022--23,2023--24,2024--25,2025--26},
  axis x line=bottom, axis y line=left, axis line style={gray!70},
  ymajorgrids, grid style={gray!20},
  tick label style={font=\footnotesize}, label style={font=\footnotesize},
  xlabel={SSHRC results year}, ylabel={Spread of committee success rates}, ylabel style={align=center},
  clip=false,
  ytick distance=0.25, yticklabel style={/pgf/number format/.cd, fixed, fixed zerofill, precision=2},
]
\draw[gray, dashed, line width=0.7pt] (axis cs:2020,0) -- (axis cs:2020,1.75);
\addplot[gray, line width=0.6pt, densely dotted, forget plot, domain=2017.6:2025.4, samples=2] {1};
\node[font=\footnotesize, text=gray!80!black, anchor=south east] at (axis cs:2025.4,1) {chance};
\addplot[color=sshrcB, line width=1.2pt, mark=square*, mark size=1.8pt] coordinates {(2018.00,0.34) (2019.00,0.31) (2020.00,0.08) (2021.00,0.10) (2022.00,0.07) (2023.00,0.06) (2024.00,0.11) (2025.00,0.03)};
\end{axis}
\end{tikzpicture}

% \caption{\footnotesize Group 1, Stream~B}
% \label{fig:sshrc_spread_common_B}
% \end{figure}


% Spread figure, dropped for the moment:
% \begin{figure}
% \centering
% % GENERATED BY sshrc_figures.py - edit that file, not this one.
% The document that inputs this file supplies the float, the caption and the label.
% Stream A. Committees: Medieval, classics and religious studies, Economics, Business (14A, 14B, 14C), Multidisciplinary humanities, Multidisciplinary social sciences (23, 23A, 23B).
% spread (Pearson chi2/df) and its whole-award floor: 2018: 1.206 (0.009); 2019: 0.110 (0.022); 2020: 0.018 (0.018); 2021: 0.045 (0.005); 2022: 0.063 (0.022); 2023: 0.158 (0.032); 2024: 0.043 (0.027); 2025: 0.126 (0.017)
\definecolor{sshrcA}{HTML}{2A78D6}
\definecolor{sshrcB}{HTML}{EB6834}
\begin{tikzpicture}
\begin{axis}[
  width=15cm, height=5.5cm,
  xmin=2017.6, xmax=2025.4, ymin=0, ymax=1.25,
  xtick={2018,2019,2020,2021,2022,2023,2024,2025}, xticklabels={2018--19,2019--20,2020--21,2021--22,2022--23,2023--24,2024--25,2025--26},
  axis x line=bottom, axis y line=left, axis line style={gray!70},
  ymajorgrids, grid style={gray!20},
  tick label style={font=\footnotesize}, label style={font=\footnotesize},
  xlabel={SSHRC results year}, ylabel={Spread of committee success rates}, ylabel style={align=center},
  clip=false,
  ytick distance=0.25, yticklabel style={/pgf/number format/.cd, fixed, fixed zerofill, precision=2},
]
\draw[gray, dashed, line width=0.7pt] (axis cs:2020,0) -- (axis cs:2020,1.25);
\addplot[gray, line width=0.6pt, densely dotted, forget plot, domain=2017.6:2025.4, samples=2] {1};
\node[font=\footnotesize, text=gray!80!black, anchor=south east] at (axis cs:2025.4,1) {chance};
\addplot[color=sshrcA, line width=1.2pt, mark=*, mark size=1.8pt] coordinates {(2018.00,1.21) (2019.00,0.11) (2020.00,0.02) (2021.00,0.05) (2022.00,0.06) (2023.00,0.16) (2024.00,0.04) (2025.00,0.13)};
\end{axis}
\end{tikzpicture}

% \caption{\footnotesize Group 2, Stream~A}
% \label{fig:sshrc_spread_above_A}
% \end{figure}



% Spread figure, dropped for the moment:
% \begin{figure}
% \centering
% % GENERATED BY sshrc_figures.py - edit that file, not this one.
% The document that inputs this file supplies the float, the caption and the label.
% Stream B. Committees: Medieval, classics and religious studies, Economics, Business (14A, 14B, 14C), Multidisciplinary humanities, Multidisciplinary social sciences (23, 23A, 23B).
% spread (Pearson chi2/df) and its whole-award floor: 2018: 0.543 (0.018); 2019: 0.096 (0.034); 2020: 0.088 (0.005); 2021: 0.096 (0.007); 2022: 0.038 (0.018); 2023: 0.038 (0.022); 2024: 0.042 (0.018); 2025: 0.107 (0.013)
\definecolor{sshrcA}{HTML}{2A78D6}
\definecolor{sshrcB}{HTML}{EB6834}
\begin{tikzpicture}
\begin{axis}[
  width=15cm, height=5.5cm,
  xmin=2017.6, xmax=2025.4, ymin=0, ymax=1.25,
  xtick={2018,2019,2020,2021,2022,2023,2024,2025}, xticklabels={2018--19,2019--20,2020--21,2021--22,2022--23,2023--24,2024--25,2025--26},
  axis x line=bottom, axis y line=left, axis line style={gray!70},
  ymajorgrids, grid style={gray!20},
  tick label style={font=\footnotesize}, label style={font=\footnotesize},
  xlabel={SSHRC results year}, ylabel={Spread of committee success rates}, ylabel style={align=center},
  clip=false,
  ytick distance=0.25, yticklabel style={/pgf/number format/.cd, fixed, fixed zerofill, precision=2},
]
\draw[gray, dashed, line width=0.7pt] (axis cs:2020,0) -- (axis cs:2020,1.25);
\addplot[gray, line width=0.6pt, densely dotted, forget plot, domain=2017.6:2025.4, samples=2] {1};
\node[font=\footnotesize, text=gray!80!black, anchor=south east] at (axis cs:2025.4,1) {chance};
\addplot[color=sshrcB, line width=1.2pt, mark=square*, mark size=1.8pt] coordinates {(2018.00,0.54) (2019.00,0.10) (2020.00,0.09) (2021.00,0.10) (2022.00,0.04) (2023.00,0.04) (2024.00,0.04) (2025.00,0.11)};
\end{axis}
\end{tikzpicture}

% \caption{\footnotesize Group 2, Stream~B}
% \label{fig:sshrc_spread_above_B}
% \end{figure}

\newpage
\section{Another way to see within stream success rate equalization}

\noindent Let $a^P_{c,t}$ and $s^P_{c,t}$ be the number of eligible and succesful applicants, respectively,s in commitee $c$, year $t$, stream $P\in\{A,B\}$. The pooled success rate in group $g$, year $t$, stream $P$ is given by:
\begin{align*}
    \rho^P_{g,t} = \dfrac{\sum_{c\in \text{group $g$}}s^P_{c,t}}{\sum_{c\in \text{group $g$}}a^P_{c,t}}
\end{align*}
Let $\tilde{s}^{P}_{c,t} = \rho_t^P\cdot a_{c,t}^P$ be the number of succesful applicants in the counterfactual where the commitee had a success rate equal to the pooled rate $\rho_{g,t}^P$. For each commitee and year, the figures below plot the difference $s_{c,t}^P - \tilde{s}_{c,t}^{P}$. If the difference lies between the dashed horizontal lines, then the error relative to the counterfactual is at most 1 award and this is observational evidence of success rate equalization within each group $\times$ stream. 

\begin{figure}
\centering
% GENERATED BY sshrc_figures.py - edit that file, not this one.
% The document that inputs this file supplies the float, the caption and the label.
% Stream A. Committees, left to right within each year: Philosophy, History, Fine arts, Literature, Sociology, Geography, Psychology (10A, 10B), Political science, Education (12A, 12B, 12C), Anthropology, Linguistics, Communications, Law and criminology.
\definecolor{sshrcA}{HTML}{2A78D6}
\definecolor{sshrcB}{HTML}{EB6834}
\begin{tikzpicture}
\begin{axis}[
  width=15cm, height=5.6cm,
  xmin=2017.6, xmax=2025.4, ymin=-5, ymax=5,
  xtick={2018,2019,2020,2021,2022,2023,2024,2025}, xticklabels={2018--19,2019--20,2020--21,2021--22,2022--23,2023--24,2024--25,2025--26},
  axis x line=bottom, axis y line=left, axis line style={gray!70},
  ymajorgrids, grid style={gray!20},
  tick label style={font=\footnotesize}, label style={font=\footnotesize},
  xlabel={SSHRC results year}, ylabel={$s^{A}_{c,t}-\tilde{s}^{A}_{c,t}$}, ylabel style={align=center},
  clip=false,
  ytick distance=1,
]
\fill[gray!12] (axis cs:2017.6,-1) rectangle (axis cs:2025.4,1);
\draw[gray, dashed, line width=0.7pt] (axis cs:2020,-5) -- (axis cs:2020,5);
\addplot[color=sshrcA, line width=1.2pt, forget plot, domain=2017.6:2025.4, samples=2] {0};
\addplot[gray!70!black, line width=0.7pt, dashed, forget plot, domain=2017.6:2025.4, samples=2] {-1};
\addplot[gray!70!black, line width=0.7pt, dashed, forget plot, domain=2017.6:2025.4, samples=2] {1};
\addplot[only marks, mark=*, mark size=1.8pt, color=sshrcA,
  mark options={fill opacity=0.6, draw opacity=0}]
  coordinates {(2018.18,0.48) (2018.26,-2.53) (2018.06,0.47) (2018.10,-0.02) (2018.14,0.47) (2017.78,1.47) (2017.90,0.98) (2017.74,-4.06) (2018.30,-0.54) (2018.22,-0.02) (2017.82,2.96) (2017.70,-1.04) (2018.02,2.46) (2017.94,-3.53) (2017.98,0.97) (2017.86,1.46) (2019.18,0.63) (2019.26,-0.12) (2019.06,-0.55) (2019.10,-1.01) (2018.78,0.42) (2018.90,-0.23) (2018.74,0.11) (2019.30,-0.15) (2019.22,0.23) (2018.82,1.65) (2018.70,0.45) (2019.02,0.51) (2018.94,-1.75) (2018.98,-0.23) (2018.86,0.05) (2019.70,0.83) (2019.74,-0.52) (2019.78,-0.24) (2019.82,0.04) (2019.86,-0.59) (2019.90,-0.24) (2019.94,-0.14) (2019.98,0.30) (2020.02,0.40) (2020.06,-0.25) (2020.10,-0.59) (2020.18,0.33) (2020.22,0.31) (2020.26,0.41) (2020.30,-0.04) (2020.70,0.30) (2020.74,0.08) (2020.78,0.08) (2020.82,-0.81) (2020.86,-0.24) (2020.90,-0.01) (2020.94,-0.13) (2020.98,-0.02) (2021.02,-0.25) (2021.06,0.10) (2021.10,-0.35) (2021.14,0.21) (2021.18,-0.34) (2021.22,-0.12) (2021.26,0.20) (2021.30,1.31) (2021.70,-0.17) (2021.74,0.64) (2021.78,0.19) (2021.82,-0.26) (2021.86,0.64) (2021.90,0.18) (2021.94,-0.09) (2021.98,0.00) (2022.02,0.46) (2022.06,0.00) (2022.10,-0.09) (2022.14,0.46) (2022.18,-0.36) (2022.22,-0.63) (2022.26,-0.36) (2022.30,-0.63) (2022.70,-0.80) (2022.74,0.18) (2022.78,2.11) (2022.82,-0.29) (2022.86,0.57) (2022.90,-0.42) (2022.94,-0.40) (2022.98,-0.37) (2023.02,-0.82) (2023.06,0.12) (2023.10,0.11) (2023.14,0.11) (2023.18,-0.48) (2023.22,0.05) (2023.26,0.18) (2023.30,0.15) (2023.70,0.31) (2023.74,-0.79) (2023.78,-0.61) (2023.82,0.25) (2023.86,-0.09) (2023.90,-0.56) (2023.94,0.30) (2023.98,0.53) (2024.02,-0.41) (2024.06,0.30) (2024.10,-0.18) (2024.14,-0.32) (2024.18,0.20) (2024.22,-0.42) (2024.26,0.73) (2024.30,0.77) (2024.70,-0.57) (2024.74,-0.29) (2024.78,-0.43) (2024.82,-0.43) (2024.86,0.00) (2024.90,0.43) (2024.94,0.14) (2024.98,0.29) (2025.02,0.14) (2025.06,0.14) (2025.10,0.57) (2025.14,0.43) (2025.18,-0.14) (2025.22,0.00) (2025.26,-0.14) (2025.30,-0.14)};
\node[font=\scriptsize, anchor=base east, text=gray!80!black] at ([yshift=12pt]axis cs:2017.6,5) {pooled rate};
\node[font=\scriptsize, anchor=base] at ([yshift=12pt]axis cs:2018,5) {50.1\%};
\node[font=\scriptsize, anchor=base] at ([yshift=12pt]axis cs:2019,5) {48.6\%};
\node[font=\scriptsize, anchor=base] at ([yshift=12pt]axis cs:2020,5) {44.6\%};
\node[font=\scriptsize, anchor=base] at ([yshift=12pt]axis cs:2021,5) {55.7\%};
\node[font=\scriptsize, anchor=base] at ([yshift=12pt]axis cs:2022,5) {54.5\%};
\node[font=\scriptsize, anchor=base] at ([yshift=12pt]axis cs:2023,5) {49.2\%};
\node[font=\scriptsize, anchor=base] at ([yshift=12pt]axis cs:2024,5) {38.0\%};
\node[font=\scriptsize, anchor=base] at ([yshift=12pt]axis cs:2025,5) {42.9\%};
\node[font=\scriptsize, anchor=base east, text=gray!80!black] at ([yshift=3pt]axis cs:2017.6,5) {within $\pm1$ award};
\node[font=\scriptsize, anchor=base] at ([yshift=3pt]axis cs:2018,5) {8/16};
\node[font=\scriptsize, anchor=base] at ([yshift=3pt]axis cs:2019,5) {12/15};
\node[font=\scriptsize, anchor=base] at ([yshift=3pt]axis cs:2020,5) {15/15};
\node[font=\scriptsize, anchor=base] at ([yshift=3pt]axis cs:2021,5) {15/16};
\node[font=\scriptsize, anchor=base] at ([yshift=3pt]axis cs:2022,5) {16/16};
\node[font=\scriptsize, anchor=base] at ([yshift=3pt]axis cs:2023,5) {15/16};
\node[font=\scriptsize, anchor=base] at ([yshift=3pt]axis cs:2024,5) {16/16};
\node[font=\scriptsize, anchor=base] at ([yshift=3pt]axis cs:2025,5) {16/16};
\end{axis}
\end{tikzpicture}

\caption{\footnotesize Group 1, Stream~A}
\label{fig:sshrc_awards_common_A}
\end{figure}

\begin{figure}
\centering
% GENERATED BY sshrc_figures.py - edit that file, not this one.
% The document that inputs this file supplies the float, the caption and the label.
% Stream B. Committees, left to right within each year: Philosophy, History, Fine arts, Literature, Sociology, Geography, Psychology (10A, 10B), Political science, Education (12A, 12B, 12C), Anthropology, Linguistics, Communications, Law and criminology.
\definecolor{sshrcA}{HTML}{2A78D6}
\definecolor{sshrcB}{HTML}{EB6834}
\begin{tikzpicture}
\begin{axis}[
  width=15cm, height=5.6cm,
  xmin=2017.6, xmax=2025.4, ymin=-5, ymax=5,
  xtick={2018,2019,2020,2021,2022,2023,2024,2025}, xticklabels={2018--19,2019--20,2020--21,2021--22,2022--23,2023--24,2024--25,2025--26},
  axis x line=bottom, axis y line=left, axis line style={gray!70},
  ymajorgrids, grid style={gray!20},
  tick label style={font=\footnotesize}, label style={font=\footnotesize},
  xlabel={SSHRC results year}, ylabel={$s^{B}_{c,t}-\tilde{s}^{B}_{c,t}$}, ylabel style={align=center},
  clip=false,
  ytick distance=1,
]
\fill[gray!12] (axis cs:2017.6,-1) rectangle (axis cs:2025.4,1);
\draw[gray, dashed, line width=0.7pt] (axis cs:2020,-5) -- (axis cs:2020,5);
\addplot[color=sshrcB, line width=1.2pt, forget plot, domain=2017.6:2025.4, samples=2] {0};
\addplot[gray!70!black, line width=0.7pt, dashed, forget plot, domain=2017.6:2025.4, samples=2] {-1};
\addplot[gray!70!black, line width=0.7pt, dashed, forget plot, domain=2017.6:2025.4, samples=2] {1};
\addplot[only marks, mark=square*, mark size=1.8pt, color=sshrcB,
  mark options={fill opacity=0.6, draw opacity=0}]
  coordinates {(2018.18,0.59) (2018.26,-3.26) (2018.06,0.50) (2018.10,-0.69) (2018.14,0.65) (2017.78,2.08) (2017.90,1.23) (2017.74,-2.02) (2018.30,-1.11) (2018.22,0.25) (2017.82,1.25) (2017.70,-0.39) (2018.02,2.35) (2017.94,-3.57) (2017.98,1.11) (2017.86,1.04) (2019.18,1.29) (2019.26,0.10) (2019.06,-0.96) (2019.10,-3.22) (2018.78,0.24) (2018.90,-1.74) (2018.74,-0.30) (2019.30,-0.71) (2019.22,0.74) (2018.82,1.11) (2018.70,0.18) (2019.02,4.89) (2018.94,-0.34) (2018.98,-1.14) (2018.86,-0.14) (2019.70,0.15) (2019.74,-0.54) (2019.78,-0.72) (2019.82,0.02) (2019.86,-0.79) (2019.90,-1.17) (2019.94,-0.08) (2019.98,0.56) (2020.02,-0.21) (2020.06,-0.20) (2020.10,-0.48) (2020.18,1.26) (2020.22,0.17) (2020.26,1.90) (2020.30,0.12) (2020.70,-0.46) (2020.74,0.77) (2020.78,0.14) (2020.82,-1.29) (2020.86,0.56) (2020.90,0.02) (2020.94,1.59) (2020.98,0.59) (2021.02,0.08) (2021.06,-1.95) (2021.10,-0.44) (2021.14,0.08) (2021.18,-0.15) (2021.22,-0.29) (2021.26,0.37) (2021.30,0.37) (2021.70,-0.39) (2021.74,0.24) (2021.78,0.73) (2021.82,-0.04) (2021.86,0.05) (2021.90,1.12) (2021.94,-0.74) (2021.98,-0.53) (2022.02,0.59) (2022.06,0.26) (2022.10,0.19) (2022.14,0.73) (2022.18,-1.02) (2022.22,-0.76) (2022.26,-0.20) (2022.30,-0.23) (2022.70,0.02) (2022.74,0.26) (2022.78,0.51) (2022.82,0.65) (2022.86,0.72) (2022.90,0.73) (2022.94,-0.85) (2022.98,-0.85) (2023.02,-0.89) (2023.06,-0.28) (2023.10,-0.68) (2023.14,0.32) (2023.18,-0.92) (2023.22,0.08) (2023.26,0.49) (2023.30,0.69) (2023.70,0.45) (2023.74,-0.68) (2023.78,-0.76) (2023.82,0.76) (2023.86,-0.20) (2023.90,-1.19) (2023.94,0.68) (2023.98,-0.31) (2024.02,-0.31) (2024.06,1.66) (2024.10,0.81) (2024.14,0.66) (2024.18,-0.78) (2024.22,1.05) (2024.26,-1.04) (2024.30,-0.80) (2024.70,0.11) (2024.74,0.10) (2024.78,-0.10) (2024.82,-0.25) (2024.86,-0.68) (2024.90,0.79) (2024.94,0.90) (2024.98,0.96) (2025.02,-0.46) (2025.06,0.32) (2025.10,-1.09) (2025.14,-0.68) (2025.18,-0.20) (2025.22,0.16) (2025.26,-0.26) (2025.30,0.39)};
\node[font=\scriptsize, anchor=base east, text=gray!80!black] at ([yshift=12pt]axis cs:2017.6,5) {pooled rate};
\node[font=\scriptsize, anchor=base] at ([yshift=12pt]axis cs:2018,5) {42.4\%};
\node[font=\scriptsize, anchor=base] at ([yshift=12pt]axis cs:2019,5) {40.3\%};
\node[font=\scriptsize, anchor=base] at ([yshift=12pt]axis cs:2020,5) {35.6\%};
\node[font=\scriptsize, anchor=base] at ([yshift=12pt]axis cs:2021,5) {48.6\%};
\node[font=\scriptsize, anchor=base] at ([yshift=12pt]axis cs:2022,5) {46.5\%};
\node[font=\scriptsize, anchor=base] at ([yshift=12pt]axis cs:2023,5) {39.7\%};
\node[font=\scriptsize, anchor=base] at ([yshift=12pt]axis cs:2024,5) {28.4\%};
\node[font=\scriptsize, anchor=base] at ([yshift=12pt]axis cs:2025,5) {35.4\%};
\node[font=\scriptsize, anchor=base east, text=gray!80!black] at ([yshift=3pt]axis cs:2017.6,5) {within $\pm1$ award};
\node[font=\scriptsize, anchor=base] at ([yshift=3pt]axis cs:2018,5) {6/16};
\node[font=\scriptsize, anchor=base] at ([yshift=3pt]axis cs:2019,5) {9/15};
\node[font=\scriptsize, anchor=base] at ([yshift=3pt]axis cs:2020,5) {12/15};
\node[font=\scriptsize, anchor=base] at ([yshift=3pt]axis cs:2021,5) {13/16};
\node[font=\scriptsize, anchor=base] at ([yshift=3pt]axis cs:2022,5) {14/16};
\node[font=\scriptsize, anchor=base] at ([yshift=3pt]axis cs:2023,5) {16/16};
\node[font=\scriptsize, anchor=base] at ([yshift=3pt]axis cs:2024,5) {12/16};
\node[font=\scriptsize, anchor=base] at ([yshift=3pt]axis cs:2025,5) {15/16};
\end{axis}
\end{tikzpicture}

\caption{\footnotesize Group 1, Stream~B}
\label{fig:sshrc_awards_common_B}
\end{figure}

\begin{figure}
\centering
% GENERATED BY sshrc_figures.py - edit that file, not this one.
% The document that inputs this file supplies the float, the caption and the label.
% Stream A. Committees, left to right within each year: Economics, Business (14A, 14B, 14C), Multidisciplinary humanities, Multidisciplinary social sciences (23, 23A, 23B).
\definecolor{sshrcA}{HTML}{2A78D6}
\definecolor{sshrcB}{HTML}{EB6834}
\begin{tikzpicture}
\begin{axis}[
  width=15cm, height=5.6cm,
  xmin=2017.6, xmax=2025.4, ymin=-5, ymax=5,
  xtick={2018,2019,2020,2021,2022,2023,2024,2025}, xticklabels={2018--19,2019--20,2020--21,2021--22,2022--23,2023--24,2024--25,2025--26},
  axis x line=bottom, axis y line=left, axis line style={gray!70},
  ymajorgrids, grid style={gray!20},
  tick label style={font=\footnotesize}, label style={font=\footnotesize},
  xlabel={SSHRC results year}, ylabel={$s^{A}_{c,t}-\tilde{s}^{A}_{c,t}$}, ylabel style={align=center},
  clip=false,
  ytick distance=1,
]
\fill[gray!12] (axis cs:2017.6,-1) rectangle (axis cs:2025.4,1);
\draw[gray, dashed, line width=0.7pt] (axis cs:2020,-5) -- (axis cs:2020,5);
\addplot[color=sshrcA, line width=1.2pt, forget plot, domain=2017.6:2025.4, samples=2] {0};
\addplot[gray!70!black, line width=0.7pt, dashed, forget plot, domain=2017.6:2025.4, samples=2] {-1};
\addplot[gray!70!black, line width=0.7pt, dashed, forget plot, domain=2017.6:2025.4, samples=2] {1};
\addplot[only marks, mark=*, mark size=1.8pt, color=sshrcA,
  mark options={fill opacity=0.6, draw opacity=0}]
  coordinates {(2017.93,2.66) (2017.96,3.13) (2017.90,1.65) (2018.01,-3.03) (2018.04,-4.40) (2018.93,0.92) (2018.96,0.35) (2018.90,0.52) (2019.01,-0.49) (2019.04,-1.29) (2019.90,-0.51) (2019.93,-0.04) (2019.96,0.34) (2020.01,0.05) (2020.04,0.15) (2020.90,-0.32) (2020.93,0.00) (2020.96,-0.18) (2020.99,0.82) (2021.01,-0.14) (2021.04,-0.18) (2021.90,-0.04) (2021.93,-0.57) (2021.96,-0.18) (2021.99,0.21) (2022.01,0.11) (2022.07,0.54) (2022.10,-0.07) (2022.90,-0.58) (2022.93,0.37) (2022.96,-0.11) (2022.99,-0.58) (2023.01,0.74) (2023.07,-0.68) (2023.10,0.84) (2023.90,-0.41) (2023.93,-0.43) (2023.96,-0.31) (2023.99,0.02) (2024.01,0.23) (2024.07,0.23) (2024.10,0.67) (2024.90,-0.01) (2024.93,-0.22) (2024.96,-0.74) (2024.99,-0.73) (2025.01,0.04) (2025.07,1.08) (2025.10,0.58)};
\node[font=\scriptsize, anchor=base east, text=gray!80!black] at ([yshift=12pt]axis cs:2017.6,5) {pooled rate};
\node[font=\scriptsize, anchor=base] at ([yshift=12pt]axis cs:2018,5) {54.7\%};
\node[font=\scriptsize, anchor=base] at ([yshift=12pt]axis cs:2019,5) {56.6\%};
\node[font=\scriptsize, anchor=base] at ([yshift=12pt]axis cs:2020,5) {47.6\%};
\node[font=\scriptsize, anchor=base] at ([yshift=12pt]axis cs:2021,5) {59.1\%};
\node[font=\scriptsize, anchor=base] at ([yshift=12pt]axis cs:2022,5) {60.7\%};
\node[font=\scriptsize, anchor=base] at ([yshift=12pt]axis cs:2023,5) {52.6\%};
\node[font=\scriptsize, anchor=base] at ([yshift=12pt]axis cs:2024,5) {44.4\%};
\node[font=\scriptsize, anchor=base] at ([yshift=12pt]axis cs:2025,5) {49.3\%};
\node[font=\scriptsize, anchor=base east, text=gray!80!black] at ([yshift=3pt]axis cs:2017.6,5) {within $\pm1$ award};
\node[font=\scriptsize, anchor=base] at ([yshift=3pt]axis cs:2018,5) {0/5};
\node[font=\scriptsize, anchor=base] at ([yshift=3pt]axis cs:2019,5) {4/5};
\node[font=\scriptsize, anchor=base] at ([yshift=3pt]axis cs:2020,5) {5/5};
\node[font=\scriptsize, anchor=base] at ([yshift=3pt]axis cs:2021,5) {6/6};
\node[font=\scriptsize, anchor=base] at ([yshift=3pt]axis cs:2022,5) {7/7};
\node[font=\scriptsize, anchor=base] at ([yshift=3pt]axis cs:2023,5) {7/7};
\node[font=\scriptsize, anchor=base] at ([yshift=3pt]axis cs:2024,5) {7/7};
\node[font=\scriptsize, anchor=base] at ([yshift=3pt]axis cs:2025,5) {6/7};
\end{axis}
\end{tikzpicture}

\caption{\footnotesize Group 2, Stream~A}
\label{fig:sshrc_awards_above_A}
\end{figure}

\begin{figure}
\centering
% GENERATED BY sshrc_figures.py - edit that file, not this one.
% The document that inputs this file supplies the float, the caption and the label.
% Stream B. Committees, left to right within each year: Economics, Business (14A, 14B, 14C), Multidisciplinary humanities, Multidisciplinary social sciences (23, 23A, 23B).
\definecolor{sshrcA}{HTML}{2A78D6}
\definecolor{sshrcB}{HTML}{EB6834}
\begin{tikzpicture}
\begin{axis}[
  width=15cm, height=5.6cm,
  xmin=2017.6, xmax=2025.4, ymin=-5, ymax=5,
  xtick={2018,2019,2020,2021,2022,2023,2024,2025}, xticklabels={2018--19,2019--20,2020--21,2021--22,2022--23,2023--24,2024--25,2025--26},
  axis x line=bottom, axis y line=left, axis line style={gray!70},
  ymajorgrids, grid style={gray!20},
  tick label style={font=\footnotesize}, label style={font=\footnotesize},
  xlabel={SSHRC results year}, ylabel={$s^{B}_{c,t}-\tilde{s}^{B}_{c,t}$}, ylabel style={align=center},
  clip=false,
  ytick distance=1,
]
\fill[gray!12] (axis cs:2017.6,-1) rectangle (axis cs:2025.4,1);
\draw[gray, dashed, line width=0.7pt] (axis cs:2020,-5) -- (axis cs:2020,5);
\addplot[color=sshrcB, line width=1.2pt, forget plot, domain=2017.6:2025.4, samples=2] {0};
\addplot[gray!70!black, line width=0.7pt, dashed, forget plot, domain=2017.6:2025.4, samples=2] {-1};
\addplot[gray!70!black, line width=0.7pt, dashed, forget plot, domain=2017.6:2025.4, samples=2] {1};
\addplot[only marks, mark=square*, mark size=1.8pt, color=sshrcB,
  mark options={fill opacity=0.6, draw opacity=0}]
  coordinates {(2017.93,2.36) (2017.96,2.41) (2017.90,0.54) (2018.01,-1.25) (2018.04,-4.06) (2018.93,0.76) (2018.96,0.76) (2018.90,0.33) (2019.01,-0.50) (2019.04,-1.35) (2019.90,-1.05) (2019.93,0.38) (2019.96,0.72) (2020.01,-0.97) (2020.04,0.91) (2020.90,1.04) (2020.93,0.01) (2020.96,0.04) (2020.99,1.04) (2021.01,-0.29) (2021.04,-1.84) (2021.90,-0.53) (2021.93,-0.21) (2021.96,-0.21) (2021.99,-0.16) (2022.01,-0.37) (2022.07,0.68) (2022.10,0.79) (2022.90,-0.46) (2022.93,-0.13) (2022.96,-0.46) (2022.99,-0.57) (2023.01,0.54) (2023.07,0.31) (2023.10,0.76) (2023.90,0.04) (2023.93,-0.07) (2023.96,-0.72) (2023.99,-0.72) (2024.01,0.67) (2024.07,0.39) (2024.10,0.39) (2024.90,-0.16) (2024.93,-0.60) (2024.96,-0.45) (2024.99,-1.02) (2025.01,0.42) (2025.07,0.98) (2025.10,0.84)};
\node[font=\scriptsize, anchor=base east, text=gray!80!black] at ([yshift=12pt]axis cs:2017.6,5) {pooled rate};
\node[font=\scriptsize, anchor=base] at ([yshift=12pt]axis cs:2018,5) {45.9\%};
\node[font=\scriptsize, anchor=base] at ([yshift=12pt]axis cs:2019,5) {46.4\%};
\node[font=\scriptsize, anchor=base] at ([yshift=12pt]axis cs:2020,5) {41.5\%};
\node[font=\scriptsize, anchor=base] at ([yshift=12pt]axis cs:2021,5) {51.8\%};
\node[font=\scriptsize, anchor=base] at ([yshift=12pt]axis cs:2022,5) {52.6\%};
\node[font=\scriptsize, anchor=base] at ([yshift=12pt]axis cs:2023,5) {44.5\%};
\node[font=\scriptsize, anchor=base] at ([yshift=12pt]axis cs:2024,5) {35.2\%};
\node[font=\scriptsize, anchor=base] at ([yshift=12pt]axis cs:2025,5) {42.9\%};
\node[font=\scriptsize, anchor=base east, text=gray!80!black] at ([yshift=3pt]axis cs:2017.6,5) {within $\pm1$ award};
\node[font=\scriptsize, anchor=base] at ([yshift=3pt]axis cs:2018,5) {1/5};
\node[font=\scriptsize, anchor=base] at ([yshift=3pt]axis cs:2019,5) {4/5};
\node[font=\scriptsize, anchor=base] at ([yshift=3pt]axis cs:2020,5) {4/5};
\node[font=\scriptsize, anchor=base] at ([yshift=3pt]axis cs:2021,5) {3/6};
\node[font=\scriptsize, anchor=base] at ([yshift=3pt]axis cs:2022,5) {7/7};
\node[font=\scriptsize, anchor=base] at ([yshift=3pt]axis cs:2023,5) {7/7};
\node[font=\scriptsize, anchor=base] at ([yshift=3pt]axis cs:2024,5) {7/7};
\node[font=\scriptsize, anchor=base] at ([yshift=3pt]axis cs:2025,5) {6/7};
\end{axis}
\end{tikzpicture}

\caption{\footnotesize Group 2, Stream~B}
\label{fig:sshrc_awards_above_B}
\end{figure}



\end{document}